% Resources for the AASTeX:
%   https://journals.aas.org/aastex-package-for-manuscript-preparation
% Author guide:
%   https://journals.aas.org/aastexguide/

\documentclass[twocolumn,twocolappendix, linenumbers, twocolappendix,trackchanges]{aastex7}

% Package imports go here.


\usepackage{amsmath}
\usepackage{textcomp}
\usepackage{gensymb}
\usepackage{graphicx}
\usepackage{mathtools, nccmath}
\usepackage{hyperref}
\usepackage[htt]{hyphenat} % Hyphenate texttt blocks
\usepackage{xspace} % grids of plots
\usepackage{subcaption}
\captionsetup{compatibility=false}

% TODO notes 
\usepackage[colorinlistoftodos]{todonotes}
\let\Oldtodo\todo
\renewcommand{\todo}[1]{\Oldtodo[inline]{#1}}
\newcommand{\eric}[2][]{\Oldtodo[inline,color=blue, #1]{Eric: #2}}
\newcommand{\bri}[2][]{\Oldtodo[inline,color=purple, #1]{Bri: #2}}
\newcommand{\TODO}[1]{\textcolor{red}{\textbf{TODO:} #1}}

% Local commands go here.
\newcommand{\docRef}{RTN-114}
\newcommand{\docUpstreamLocation}{\url{https://github.com/lsst/rtn-114}}
\newcommand{\degsq}{$\rm deg^{2}$}

%\providecommand{\secref}[1]{\hyperref[#1]{\autoref{#1}}}
\providecommand{\secref}[1]{\hyperref[#1]{\S\ref{#1}}}
\providecommand{\appref}[1]{\hyperref[#1]{Appendix \ref{#1}}}
\providecommand{\tabref}[1]{\hyperref[#1]{\autoref{#1}}}
\providecommand{\figref}[1]{\hyperref[#1]{\autoref{#1}}}
\providecommand{\eqnref}[1]{\hyperref[#1]{Eq.~\ref{#1}}}
\providecommand{\recref}[1]{\hyperref[#1]{REC-\ref{#1}}}


% All graphics paths here 
\graphicspath{./}

%%%%%%%%%%%%%%%%%%%%%%%%%%%%%%%%%%%%%%%%%%%%%%%%%%%%%%%%%%%%%%%%%%%%%%%%%%%%%%%%
%%
%% The following section outlines numerous optional output that
%% can be displayed in the front matter or as running meta-data.
%%
%% Running header information. A short title on odd pages and 
%% short author list on even pages. Note that this
%% information may be modified in production.
%%\shorttitle{AASTeX v7 Sample article}
%%\shortauthors{The Terra Mater collaboration}
%%
%% Include dates for submitted, revised, and accepted.
%%\received{February 1, 2025}
%%\revised{March 1, 2025}
%%\accepted{\today}
%%
%% Indicate AAS Journal the manuscript was submitted to.
%% Note that this command adds "Submitted to " the argument.
%%\submitjournal{PSJ}
%%
%% You can add a light gray and diagonal water-mark to the first page 
%% with this command:
%% \watermark{text}
%% where "text", e.g. DRAFT, is the text to appear.  If the text is 
%% long you can control the water-mark size with:
%% \setwatermarkfontsize{dimension}
%% where dimension is any recognized LaTeX dimension, e.g. pt, in, etc.
%%%%%%%%%%%%%%%%%%%%%%%%%%%%%%%%%%%%%%%%%%%%%%%%%%%%%%%%%%%%%%%%%%%%%%%%%%%%%%%%
%%
%% Use this command to indicate a subdirectory where figures are located.
%%\graphicspath{{./}{figures/}}



\begin{document}
\doi{10.71929/rubin/3019817}
\input{authors}
\correspondingauthor{Eric C. Bellm}

\date{\today}
\title{Alert Production with the Vera C. Rubin Observatory}

% This can write metadata into the PDF.
% Update keywords and author information as necessary.
\hypersetup{pdftitle={\@title}, pdfauthor={\@author}, pdfkeywords={\@keywords}}

\input{abstract}

%% Keywords should appear after the \end{abstract} command.
%% The AAS Journals now uses Unified Astronomy Thesaurus concepts:
%% https://astrothesaurus.org
%% You can use the \uat command to link your UAT concepts back its source.
%% \keywords{\uat{Galaxies}{573} --- \uat{Cosmology}{343}}

\listoftodos

\input{intro}

\section{Pipeline}

%\note{following LDM-151}
\todo{consider a (simplified?) pipeline graph}

\subsection{Template Generation}

%\note{Pre-DR1 this may require description independent of DRP.}
%
%High-level discussion of DCR correction could go here.

\subsection{Preload}

\subsection{Single Frame Processing}

%\note{Construction of calibration products assumed described in \cite{PSTN-026}.}

ISR, PSF and Background fitting, photometric and astrometric calibration.  Likely major commonalities with DRP.

\subsection{Image Differencing}

\subsection{Source Detection and Measurement}

Including discussion of point source, dipole, and streak detection on difference images.  

%\note{Algorithmic basics could be discussed elsewhere}

\subsection{Initial Filtering} \label{sec:filtering}

\texttt{FilterDiaSourceCatalogTask} in \texttt{lsst.ap.association}

\subsection{Reliability scoring}

Image-differencing searches for transients typically contend with high rates of false positives.
Among the raw detections, artifacts may dominate real astrophysical sources by an order of magnitude.
These may be due to imperfectly-corrected instrument signatures, astrometric mismatches between the template and the science image, cosmic rays, or algorithmic failures in the differencing.
Modern surveys employ machine-learned classifiers to winnow these candidates \citep[e.g.,][]{Bailey:07:NearbySNFML,Bloom:12:RealBogus,Brink:13:RB2,Wright:15:PS1ML,Goldstein:15:DESRealBogus,Duev:19:Braii}.
These have steadily improved in performance and sophistication as the classifiers transitioned from Random Forest models trained on manually-constructed feature sets to deep neural networks trained directly on image pixels.
Some approaches have combined both detection and scoring without using the difference image \citep{Sedaghat:17:Transinet,Acero-Cuellar:22:NoDifference}.

\todo{overview of the algorithm, training, and performance of the ML spuriousness score; detailed discussion likely deferred to separate paper}

\subsection{Catalog Transformation}

\texttt{TransformDiaSourceCatalogTask} in  \texttt{lsst.ap.association}

\eric{schema tables?}

\subsection{Source Association}

Due to AP's real-time nature, it is not possible to perform \textit{post-facto} source association as in the annual data releases.
Instead, spatial association is performed on-the-fly by the \texttt{DiaPipelineTask} within the \texttt{lsst.ap.association} package.
A dedicated Alert Production Database (APDB; \S \ref{sec:db}) holds the current state of the system.

During the preload step (\S \ref{sec:preload}), DIAObjects and SSObjects overlapping the expected field of view of the image are stored in the local prompt processing worker butler repository.
The new DIASources are first spatially associated with the existing DIAObjects by finding the closest match within a maximum distance of one arcsecond.
Pairs with the closest spatial separations are joined first to minimize misassociations.
When a match is found, the DIASource is added to the DIAObject's list of measurements.
Unassociated DIASources are then spatially associated with the predicted positions of the SSObjects at the time of the exposure.
For DIASources with no matching DIAObject or SSObject, a new DIAObject is created and the DIASource is added to it.



%\note{does \cite{PSTN-045} describe the APDB/PPDB implementation?}

\subsection{Alert Generation}

Each DIASource not filtered for quality purposes (\S \ref{sec:filtering}) produces an alert.
To enable downstream users to quickly identify events of interest, the alert packets contain a range of information.  
They include:

\begin{itemize}
    \item The triggering DIASource record
    \item The corresponding DIAObject or SSObject record
    \item Any DIASource and DIAForcedSource records that exist, and difference image noise estimates where they do not, taken from the previous 12 months.
    \item Cutout images of the science, template, and difference images.
\end{itemize}
\eric{confirm contents}

The cutout images are packaged in the FITS format using the astropy \texttt{CCDData} class and include mask and variance planes, an estimate of the PSF, a World Coordinate System, and other metadata.
Most cutouts are 30$\times$30 pixels but are expanded to a larger size for extended sources.
The maximimum cutout size of \todo{confirm} allows cutouts to contain trailed Near Earth Asteroids traveling slower than ten degrees per day.

Alerts are transmitted and stored in Apache Avro\footnote{\url{https://avro.apache.org/}} format, a strongly-schemaed compact binary serialization.
However, the Avro standard includes its schemas as uncompressed ASCII.
To conserve bandwidth, we therefore send alerts without including the schemas required to deserialize them.
Instead, the schemas are stored in a separate schema registry and are referenced by a unique identifier in the alert packet.

\eric{sizing estimates here or in initial performance section?}

\subsection{Alert Distribution} \label{sec:alertdist}

We use the Kafka distributed streaming platform\footnote{\url{https://kafka.apache.org/}} to transmit alerts to brokers.
Kafka provides a fault-tolerant, low-overhead mechanism to transmit large volumes of data to multiple consumers and is widely used both in industry and in astronomy. 
\eric{cite patterson, GCN, Hopskotch, other Rubin usage...}

The Alert Distribution Kafka system is deployed at the USDF using Kubernetes.
\bri{details of the deployment}
\eric{schema management}

Because of the large size of the alert packets and the short latency requirement, bandwidth out of the US Data Facility (USDF) at SLAC limits the number of direct recepients of the full alert stream.
Seven community alert brokers were selected through a proposal process \eric{cite} to receive the full alert stream directly from the USDF:
\eric{list}
Two additional brokers will operate downstream of a full-stream broker.


\subsection{Alert Filtering Service}

ANTARES system

community alert filters

\subsection{Forced Photometry}

\subsection{Source Injection}

\subsection{Metrics and Monitoring}
\input{data}
\section{Processing Environment}


\subsection{Prompt Processing Framework} \label{sec:pp}

\subsection{Databases} \label{sec:db}

\subsection{Alert Archive}

\subsection{Metrics and Dashboards}

\subsection{Catchup Processing}
\section{Initial Performance} \label{sec:performance}

\subsection{Template Properties}

\subsection{Alert Rate} \label{sec:rates}

\subsection{Alert Purity} \label{sec:purity}

Raw and after ML scoring


\subsection{Alert Completeness}

\subsection{Alert Latency}

\subsection{Photometric Precision}

\subsection{Astrometric Precision}

\subsection{Association Accuracy}

\subsection{Association of Known Solar System Objects}

\jake{completeness, astrometric accuracy}

\subsection{Direct Source Catalogs}

\subsection{Crowded Field Performance}
%\section{Discussion}
\section{Conclusion} \label{sec:conclusion}

\eric{summarize}




%% Modify acknowledgments as needed.
\begin{acknowledgments}
This material is based upon work supported in part by the National Science Foundation through Cooperative Agreements AST-1258333 and AST-2241526 and Cooperative Support Agreements AST-1202910 and AST-2211468 managed by the Association of Universities for Research in Astronomy (AURA), and the Department of Energy under Contract No.\ DE-AC02-76SF00515 with the SLAC National Accelerator Laboratory managed by Stanford University.
Additional Rubin Observatory funding comes from private donations, grants to universities, and in-kind support from LSST-DA Institutional Members.
\end{acknowledgments}

%% To help institutions obtain information on the effectiveness of their
%% telescopes the AAS Journals has created a group of keywords for telescope
%% facilities.
%
%% Following the acknowledgments section, use the following syntax and the
%% \facility{} or \facilities{} macros to list the keywords of facilities used
%% in the research for the paper.  Each keyword is check against the master
%% list during copy editing.  Individual instruments can be provided in
%% parentheses, after the keyword, but they are not verified.

\vspace{5mm}
\facility{Rubin:Simonyi (LSSTCam), Rubin:USDAC}

%% Similar to \facility{}, there is the optional \software command to allow
%% authors a place to specify which programs were used during the creation of
%% the manuscript. Authors should list each code and include either a
%% citation or url to the code inside ()s when available.
\software{Rubin Data Butler \citep{2022SPIE12189E..11J},
          LSST Science Pipelines \citep{PSTN-019},
          LSST Feature Based Scheduler v3.0  \citep{peter_yoachim_2024_13985198, Naghib_2019}
          Astropy \citep{2013A&A...558A..33A, 2018AJ....156..123A, 2022ApJ...935..167A}
          }

%% Appendix material should be preceded with a single \appendix command.
%% There should be a \section command for each appendix. Mark appendix
%% subsections with the same markup you use in the main body of the paper.

%% Each Appendix (indicated with \section) will be lettered A, B, C, etc.
%% The equation counter will reset when it encounters the \appendix
%% command and will number appendix equations (A1), (A2), etc. The
%% Figure and Table counter will not reset.

%% \appendix

%\printglossaries

% Include all the relevant bib files so that references can be found from
% lsst-texmf.
% https://lsst-texmf.lsst.io/lsstdoc.html#bibliographies
\bibliographystyle{aasjournal}
\bibliography{local,lsst,lsst-dm,refs_ads,refs,books}

\end{document}

